\section{Summary}
\label{sec:conclusions}

The Goldfarb--Idnani method is often presented as a sequence of updates, which
makes it look more intricate than it is. Read from the invariants it is short. One
matrix $J$ is carried, defined by $J\T G J = I$ and $J\T A = [R;0]$: its columns
are a $G$-orthonormal basis of $\R^n$, arranged so that the trailing block spans
the working set's feasible subspace. Every quantity the iteration needs is then one
matrix--vector product or one triangular solve away, and every quantity it
\emph{consumes} has a closed form in $(G, A, n_*)$ alone
(Proposition~\ref{prop:invariance}), which is why the representation may be
chosen for speed, and why two implementations holding different factors compute
the same step.

The iteration is one affine path (Proposition~\ref{prop:path}) along which
stationarity holds identically while the entering multiplier rises and the
working-set multipliers fall. Two limits cut it short, and which binds decides
whether a constraint is added or dropped. The objective advances by an exact closed
form (Proposition~\ref{prop:increment}), positive in both step directions, so the
method is a dual ascent and, absent degeneracy, finite
(Proposition~\ref{prop:ascent}). When the path is blocked in both directions the
vector $r$ is a Farkas certificate (Proposition~\ref{prop:farkas}): the
infeasibility verdict is a proof, and the solver already holds it.

Two properties of the classical method are easy to overlook, and the extensions
of Sections~\ref{sec:pdas} and~\ref{sec:sweep} trade one away and exploit the
other. The first is that linear dependence among the working-set normals is
unreachable: the step rules exclude it (Proposition~\ref{prop:rank}), so no rank
test is needed anywhere. Guessing the active set forfeits it, and a primal--dual
active-set fast path therefore admits no finite-termination theorem for general
constraints: there is no $P$-matrix structure to appeal to once the working-set
system stops being a principal submatrix. The second is that the factors depend
on $G$ and the working set but not on the linear term, so a family of programs
differing only in $a$ is solvable from one factorisation, and a stale active set
repairable rather than disposable.

Both extensions are unguaranteed and both are safe, for the same reason: strict
convexity makes the KKT conditions sufficient (Proposition~\ref{prop:sufficient}),
so a candidate can be certified instead of trusted. The strict convexity that
buys it is a choice of scope rather than a boundary of the approach: the dual
method extends to semidefinite $G$, with finite termination, by
Boland~\cite{boland1997}, and by proximal regularisation in the solver
of~\cite{arnstrom2022daqp}. Of the ideas here, that one travels furthest. A method that may fail, but that can recognise its own success,
needs a certificate and not a convergence theory.

One observation generalises past this solver. Below a few hundred variables a
solve costs what its array operations cost to dispatch, not what its
arithmetic costs, so halving the flop count buys nothing and halving the
\emph{iteration} count buys everything: the fast path of Section~\ref{sec:pdas}
gains about as much over the exact walk as rewriting that walk in another language
does, and gains most on the family where the rewrite gains least
(Section~\ref{ssec:compare}). Effort spent on kernels is
wasted where the bottleneck is the interpreter, and effort spent on iteration
counts is wasted where it is not; knowing which regime one is in is the whole of
the decision.
